\color{unidarkgreen}{\subsection{Контроль отсутствия напряжения}}
\color{black}

\begin{enumerate}[label=\arabic{section}.\arabic{subsection}.\arabic*,  align=left, labelsep=4pt, leftmargin=0pt, itemindent=57pt]

\item
Ввод в работу ЗНР осуществляется с помощью уставки \textbf{«Работа КОН».}

\item
Для исключения пуска пожаротушения не отключенного от сети ШР обеспечиваются следующие виды контроля:
\begin{itemize}
\item отсутствие тока на стороне ВВ, контролируется порогом срабатывания \textbf{«Iср КОН ВВ ЗНР»};
\item отсутствие напряжения на стороне, контролируется порогом срабатывания \textbf{«Uср КОН НН»}.
\end{itemize}

\begin{figure}[H]
\centering
\includegraphics[width=0.7\textwidth,height=0.7\textheight,keepaspectratio]{\fbpath/161.1201 КОН/img/kon.pdf}
\captionof{figure}{Функциональный блок КОН ШР}\label{pict_kon:KON}
\end{figure}

\small
\begin{longtable}{|>{\centering\arraybackslash}m{5.3cm}|>{\centering\arraybackslash}m{4.2cm}|>{\centering\arraybackslash}m{1cm}|>{\centering\arraybackslash}m{3.3cm}|>{\centering\arraybackslash}m{1.8cm}|}
% Первая страница: полный заголовок
\caption{Перечень уставочных параметров КОН ШР\hfill\vspace{-0.5\baselineskip}}\label{kon:tbl1}\\
\hline
\rowcolor{gray!30}
Наименование уставки & Диапазон & Шаг & Значение по умолчанию & Единица измерения \\
\hline
\endfirsthead

% Промежуточные страницы: «Продолжение таблицы»
\caption*{\hspace{3pt}\emph{Продолжение таблицы \ref{kon:tbl1}}\hfill\ }\vspace{-0.5\baselineskip} \\
\hline
\rowcolor{gray!30}
Наименование уставки & Диапазон & Шаг & Значение по умолчанию & Единица измерения \\
\hline
\endhead

%===m>КОН|-
%===m
\end{longtable}
\normalsize

\end{enumerate}